\documentclass[12pt,a4paper]{article}
\usepackage[utf8]{inputenc}
\usepackage[T1]{fontenc}
\usepackage{geometry}
\usepackage{graphicx}
\usepackage{blindtext}
\usepackage{floatflt}
\newcommand{\changefont}[3]{\fontfamily{#1} \fontseries{#2}
                            \fontshape{#3} \selectfont}

\geometry{a4paper,left=20mm,right=25mm, top=12mm, bottom=12mm}
\parindent 0px 
\parskip 0px
\author{Tim}

\begin{document}
	\thispagestyle{empty}
	\changefont{pag}{m}{n}
	\begin{figure}[h!]
		\centering
		\includegraphics[width=0.8\textwidth]{fsilogo.eps}
	\end{figure}
Wir sind eine Gruppe Studenten der Studiengänge Informatik, Bioinformatik, \\ Medieninformatik, Medizininformatik und Kognitionswissenschaften und kümmern uns um personelle, politische sowie soziale Themen des Studiums hier am Wilhelm-Schickard-Institut. \\
% too much?
\begin{itemize}
	\item \textbf{Neu in Tübingen?} Wir betreuen jedes Semester neue Studienanfänger mit zahlreichen Angeboten und erleichtern somit das Kennenlernen und den Studienanfang.
	\item \textbf{Wir organisieren Feste und Feiern!} Jedes Jahr findet hier auf dem Sand das große Sommerfest statt, ein Mal pro Semester veranstalten wir zusammen mit der Fachschaft Psychologie unser legendäres Clubhausfest.
	\item Ihr habt eine Idee oder euch ist etwas aufgefallen, was für viele Studenten von Interesse ist? \textbf{Sprecht uns an oder schreibt uns eine Mail!}
	\item \textbf{Wir bieten euch} Services wie unser Prüfungsprotokolle-Interface und diverse Mailinglisten an, die euch die Studiumsorganisation und den Kontakt mit Kommilitonen erleichtern. \\ Weitere Infos dazu findet ihr auf unserer Homepage.
	\item \textbf{Wir machen Gremienarbeit.} Wir vertreten also die Interessen aller Informatikstudenten in Gremien wie dem Fakultätsrat, der Studienkommission und anderen Universitätsgremien und setzen uns mit unserem Mitspracherecht für die Belange der Studenten ein.
\end{itemize}

\textbf{Interesse?}
Neue Mitglieder sind in der Fachschaft stets willkommen! Schau einfach hier im Fachschaftszimmer vorbei oder komm zu unserer nächsten Sitzung.
\textbf{Den aktuellen Sitzungstermin findest du auf unserer Webseite.} \\
\vspace*{1cm} 

\textbf{Hinweise zum Fachschaftszimmer}
\begin{itemize}
	\item Dieser Raum ist kein studentischer Arbeitsraum. Wir werfen zwar niemanden grundlos raus, meist herrschen jedoch keine geeigneten Bedingungen zum Lernen oder zum Bearbeiten von Übungsblättern. Platz zum Arbeiten findet ihr (fast) immer in den Räumen \textbf{C119}, \textbf{B033} und \textbf{A328.}
	\item Wir verkaufen über unser \textbf{Verkaufssystem} Süßigkeiten und gekühlte Getränke. \\ Wenn du eine Flasche aus dem Kühlschrank nimmst, lege auch bitte wieder eine (oder gleich mehrere) nach.
\end{itemize}

\vfill
\begin{flushright}
	\includegraphics[width=0.2\textwidth]{qrcode.png}
\end{flushright}
\end{document}
